\chapter*{Tóm Tắt Đồ Án}
\addcontentsline{toc}{chapter}{Tóm Tắt Đồ Án}

\section*{Tóm tắt tiếng Việt}
Ngành du lịch tự túc tại Miền Trung Việt Nam đang phát triển mạnh mẽ, tuy nhiên du khách thường đối mặt với các khó khăn lớn: khó sắp xếp thứ tự các điểm tham quan sao cho hợp lý, thiếu địa chỉ chính xác của các quán ăn đặc sản bản địa, lịch trình cứng nhắc không thích ứng khi thời tiết xấu, và khó khăn trong việc phân bổ chi phí khi đi theo nhóm.

Đồ án \textbf{Travel Trips Central VietNam} giải quyết toàn diện các vấn đề trên thông qua việc xây dựng một ứng dụng du lịch thông minh đa nền tảng (\en{Mobile \& PC Desktop App-First}) ứng dụng \term{Trí tuệ nhân tạo tạo sinh} (\en{Generative AI}) kết hợp cơ sở dữ liệu lớn được thu thập và làm giàu tự động bởi động cơ \term{AI Autonomous Web Crawler} về 11 tỉnh/thành phố Miền Trung -- Tây Nguyên sau sáp nhập. Các đóng góp chính của đề tài bao gồm:
\begin{itemize}
    \item \textbf{Động cơ AI Autonomous Data Crawler:} Tự động tìm kiếm, bóc tách và nạp sâu rộng hơn 250+ địa danh thắng cảnh, di tích lịch sử và ẩm thực đặc sản bản địa chính gốc cho 11 tỉnh/thành phố vào cơ sở dữ liệu MongoDB Atlas.
    \item \textbf{Động cơ Lập lịch trình thông minh không trùng lặp:} Ứng dụng mô hình \en{Google Gemini 2.5 Flash} kết hợp giải thuật đồ thị đỉnh đã duyệt (\en{Visited Set Graph}) đảm bảo sinh lịch trình cá nhân hóa theo từng mốc thời gian (Sáng - Trưa - Chiều - Tối) với tỷ lệ trùng lặp địa điểm đạt 0\%.
    \item \textbf{Bộ tiện ích thông minh chuyên sâu:} Tiện ích chia tiền nhóm (\en{Group Bill Splitter}), Đổi lịch tránh mưa tự động (\en{Weather Smart Re-scheduler}) dựa trên dữ liệu khí tượng \en{Open-Meteo}, và Thẻ vé hành trình ngoại tuyến (\en{Offline Travel Pass}).
\end{itemize}

\textbf{Từ khóa:} \textit{Trí tuệ nhân tạo, Du lịch Miền Trung -- Tây Nguyên, 11 Tỉnh thành sáp nhập, AI Crawler, Gemini AI, Lập lịch trình tự động, Travel Trips Central VietNam, Vue.js, Node.js, MongoDB, Visited Set Graph.}

\vspace{1cm}
\section*{Abstract (English)}
Self-guided travel in Central Vietnam and Central Highlands is experiencing substantial growth. However, travelers frequently encounter critical pain points: difficulties in optimizing spatial travel routes, lack of verified addresses for authentic local eateries, rigid itineraries that fail during adverse weather conditions, and complexities in group expense allocation.

The \textbf{Travel Trips Central VietNam} project addresses these challenges by developing a cross-platform, app-first intelligent travel application powered by Generative AI (\en{Google Gemini AI}) integrated with an autonomous AI Web Crawler spanning 11 merged Central Vietnam and Central Highlands provinces. The core contributions of this study include:
\begin{itemize}
    \item \textbf{Autonomous AI Data Crawler Engine:} Automatically discovers, extracts, and enriches over 250+ verified tourist attractions, culinary specialties, and accommodations across 11 merged provinces.
    \item \textbf{Smart Non-Repeating Itinerary Engine:} Combines \en{Google Gemini 2.5 Flash} with a \en{Visited Set Graph} algorithm to generate structured, time-slotted daily itineraries with a 0\% destination duplication rate.
    \item \textbf{Advanced Group \& Adaptive Features:} Interactive Group Bill Splitter, Weather-Adaptive Re-scheduler leveraging real-time \en{Open-Meteo} telemetry, and an Offline Travel Pass.
\end{itemize}

\textbf{Keywords:} \textit{Artificial Intelligence, Central Vietnam Tourism, Merged Provinces, AI Web Crawler, Travel Trips Central VietNam, Gemini AI, Automated Itinerary Planning, Vue.js, Node.js, MongoDB, Visited Set Graph.}
